% SOURCE_AUTHORITY: sga5_fr_workpass.tex, lines 7268--7291 (LF-normalized unit).
% SOURCE_SHA256: 791F4EFFC5E02832D5D77ED03518C8156D6F07E4C8238B03545DB93D883FBB28
\begin{statement}{Definição 4.1.2}
Seja $X$ um sistema projetivo indexado por $\ZZ$. Diz-se que $X$ é essencialmente constante se existe um inteiro $p$ tal que, para todos os pares $(r,s)$ de inteiros maiores ou iguais a $p$ e que verificam a desigualdade $s\ge r$, o morfismo de transição
\[
X_s\longrightarrow X_r
\]
seja um isomorfismo. É claro que basta exigir quando $r=p$.
\end{statement}

\begin{statement}{Proposição 4.1.3 (Propriedade (MLAR) e graduados)}
As seguintes afirmações para um objeto $X$ da categoria $P$ são equivalentes:
\begin{enumerate}[label=(\roman*)]
\item $X$ verifica (MLAR).
\item O sistema projetivo com valores na categoria dos objetos de $\cC$ graduados por $\NN$
\[
r\longmapsto (\gr^\alpha(X_{\alpha+r}))_{\alpha\in\NN},
\]
onde os morfismos de transição são evidentes, é essencialmente constante.
\item O sistema projetivo com valores na categoria dos objetos de $\cC$ graduados por $\NN$
\[
n\longmapsto (\gr^p(X_n))_{p\in\NN},
\]
onde os morfismos de transição são evidentes, verifica (MLAR).
\end{enumerate}
\end{statement}
